\chapter{Аудит масштабного семени ньютоновской постоянной}
\label{ch:v10-induced-newton-scale-seed-audit}

\section{Контракт физического семени}

Предыдущий гейт свёл абсолютную ньютоновскую постоянную к выбору величины
$m$ размерности $L^{-2}$. Кандидат считается физическим семенем только при
одновременном выполнении шести условий: правильной размерности, внутренней
доступности, независимости от искомых $q$ и $G$, выбора общим родителем,
типизированного отображения в $(A,B)$ и разрушения абсолютной масштабной
орбиты.

Проверены десять кандидатов: кривизна роста, обратная площадь ячейки,
квадраты спектрального обрезания, часового волнового числа, теплового
КМС-волнового числа и дираковской щели, корень топологической плотности,
щель вакуумного гессиана, dimensional transmutation и обратная наблюдаемая
ньютоновская площадь.

Их матрица имеет ранг $5$, максимальный результат равен $4/6$, а вектор
полных проходов равен
\begin{equation}
 (0,0,0,0,0,0,0,0,0,0)^T.
 \label{eq:v10-newton-seed-pass-vector}
\end{equation}
Все размерно допустимые кандидаты формально вкладываются в
$A=\alpha m^2$, $B=\beta m$, но это размерное соответствие не является
механизмом выбора.

\section{Общая масштабная нулевая мода}

Относительные связи между семенем, площадью ячейки, энергией часов,
проводимостью и кривизной кодируются матрицей
\begin{equation}
 C_m=
 \begin{pmatrix}
 1&1&0&0&0\\
 1&0&-2&0&0\\
 1&0&0&-2&0\\
 1&0&0&0&-1
 \end{pmatrix}
 \label{eq:v10-newton-seed-relative-map}
\end{equation}
в логарифмических переменных $(m,q,E_C,\kappa,\Lambda)$. Точно
\begin{equation}
 \rank C_m=4,
 \qquad \operatorname{nullity}C_m=1,
 \qquad C_m(-2,2,-1,-1,-2)^T=0.
 \label{eq:v10-newton-seed-relative-rank}
\end{equation}
Поэтому все текущие внутренние связи сохраняют инварианты
\begin{equation}
 mq,\qquad \frac{m}{E_C^2},\qquad
 \frac{m}{\kappa^2},\qquad \frac{m}{\Lambda}.
 \label{eq:v10-newton-seed-invariants}
\end{equation}
Одна независимая строка, фиксирующая $m$, повысила бы ранг до $5$, но
текущий корпус такой строки не содержит.

\section{Два формальных разрушителя орбиты}

Dimensional transmutation может создать независимую величину
\begin{equation}
 m_{\rm DT}=\mu^2
 \exp\!\left(-2\int^{g(\mu)}\frac{dg}{\beta(g)}\right),
 \label{eq:v10-newton-seed-dimensional-transmutation}
\end{equation}
но Том X ещё не вывел требуемые ненулевую $\beta$-функцию, граничное
значение связи и типизированное отображение результата в коэффициенты
кривизны.

Наблюдаемая ньютоновская площадь даёт
\begin{equation}
 m=\frac{1}{16\pi\beta g_N},
 \label{eq:v10-newton-seed-observed-G}
\end{equation}
однако это лишь обращение целевого равенства $16\pi\beta g_Nm=1$.

\begin{theorem}[Относительные калибровки не выбирают семя]
Текущие связи между $m$, $q$, $E_C$, $\kappa$ и $\Lambda$ имеют ровно одну
общую масштабную нулевую моду. Ни один из десяти кандидатов не даёт
одновременно независимый, родительски выбранный и некруговой источник $m$.
\end{theorem}

\begin{nogocriterion}[Запрет ньютоновского обратного якоря]
Нельзя подставлять наблюдаемое $G$ в формулу для $m$, а затем объявлять
полученное $m$ выводом $G$. Dimensional transmutation допустима как новый
маршрут только после отдельного вывода её $\beta$-данных общим родителем.
\end{nogocriterion}

\begin{protocol}[Статус аудита масштабного семени]\begin{enumerate}\raggedright
\item проверенных кандидатов: \textbf{$10$};
\item полных проходов: \textbf{$0/10$};
\item ранг матрицы кандидатов: \textbf{$5$};
\item максимальный результат: \textbf{$4/6$};
\item относительная карта: \textbf{ранг/ядро $4/1$};
\item происхождение масштабного семени: \textbf{$0/1$};
\item абсолютная ньютоновская постоянная: \textbf{$0/1$};
\item следующий гейт: \textbf{родитель $\beta$-данных dimensional transmutation}.
\end{enumerate}\end{protocol}

Машинный сертификат сохранён в
\path{s2t_v10_cell_birth_four_volume_induced_newton_scale_seed_candidate_audit_gate_results.json}.